% Standalone source-record documentation; no research attempts.
\documentclass[11pt]{article}
\usepackage[margin=25mm]{geometry}
\usepackage{fontspec,unicode-math}
\setmainfont{Latin Modern Roman}
\setmathfont{Latin Modern Math}
\usepackage{amsmath,amssymb,mathtools,xurl,hyperref}
\hypersetup{colorlinks=true,urlcolor=blue,linkcolor=blue}
\setcounter{secnumdepth}{0}
\setlength{\parindent}{0pt}
\setlength{\parskip}{5pt}
\setlength{\emergencystretch}{3em}
\newcommand{\sref}[2]{\hyperlink{source:#1:#2}{[S#2]}}
\newcommand{\cataloguescope}[1]{\paragraph{Source record.}#1}
\newcommand{\floor}[1]{\left\lfloor#1\right\rfloor}
\begin{document}
% UnsolvedMath ID 30002239; source catalogue code OWR-12177-008
\setcounter{section}{30002238}
\section{Polynomial Freiman–Ruzsa Conjecture over Binary Vector Spaces}\label{30002239}
{\small\sffamily UnsolvedMath ID 30002239 · Computer Science}
\subsection{Definitions and mathematical statement}

\paragraph{Shared notation.}Standard number systems are $\mathbb N,\mathbb Z,\mathbb Q,\mathbb R,\mathbb C$. Specialized notation and index ranges follow the cited source and the definitions below.


In $\mathbb F_2^n$, $A+A=\{a+b:a,b\in A\}$ uses coordinatewise addition modulo two. The doubling constant is $|A+A|/|A|$; $\operatorname{span}(A)$ is the smallest linear subspace containing $A$. Polynomial losses in $K$ mean fixed powers independent of $n$. The ordinary small-doubling hypothesis assumes $A\ne\varnothing$ and $K\ge1$.
\paragraph{Record scope and qualifications.}
The catalogue title is retained as metadata. Named certificates, contracts, bounds or frameworks in that title are not asserted as results of the original source extract.
The supplied snapshot labels this record open, but its literature assessment cites the solution of the characteristic-two Marton/Polynomial Freiman-Ruzsa conjecture for the intended nonempty-set formulation. It is retained as a historical source formulation, not asserted to remain an open conjecture.

Source-record limitation: The supplied assessment flags a survey extraction with finite-field solutions; importing its binary PFR formulation as an open target would conflict with that supplied status evidence. The supplied wording does not exclude $A=\varnothing$: then $A^{\prime}=\varnothing$ and $|\operatorname{span}(A^{\prime})|=1$, so its displayed span bound cannot hold. The solved-literature comparison applies to the ordinary nonempty-set formulation; the source statement is preserved.
\paragraph{Statement recorded in the supplied source.}
If $A\subseteq\mathbb F_2^n$ satisfies $|A+A|\leq K|A|$, is there a subset $A^{\prime}\subseteq A$ with $|A^{\prime}|\geq |A|/\operatorname{poly}(K)$ and $|\operatorname{span}(A^{\prime})|\leq\operatorname{poly}(K)|A^{\prime}|$? \sref{30002239}{2}
\paragraph{Source status.}
The supplied catalogue labels this source record open; that metadata was checked against the supplied snapshot on 2026-10-11. The label is preserved for provenance and is not an independent verification of present open status or mathematical correctness. The catalogue literature assessment is dated 2026-08-21.
\subsection{Short English statement}
Relate a small-doubling binary set to a large subset with comparably small linear span.
\subsection{Sources}
\begin{itemize}
\item[\hypertarget{source:30002239:1}{S1}] Ulam AI and the UnsolvedMath Contributors, UnsolvedMath: A Curated Collection of Open Mathematics Problems, record 30002239 (OWR-12177-008). Catalogue curation and metadata: CC BY 4.0. The source extract is retained; definitions, English summaries and qualifications are editorial additions.
\url{https://huggingface.co/datasets/ulamai/UnsolvedMath}.
\item[\hypertarget{source:30002239:2}{S2}] Complexity Theory (2013), Oberwolfach Reports, DOI 10.4171/owr/2012/54, pp. 3267--3304
\url{https://doi.org/10.4171/owr/2012/54}.
\item[\hypertarget{source:30002239:3}{S3}] Gowers--Green--Manners--Tao, Marton conjecture in characteristic two
\url{https://arxiv.org/abs/2311.05762}.
\item[\hypertarget{source:30002239:4}{S4}] Kazemi Moghadam, odd-characteristic finite-field case
\url{https://arxiv.org/abs/2512.04433}.
\item Source-specific Oberwolfach Reports licences, including CC BY-SA 4.0 where applicable, are retained; see THIRD\_PARTY\_NOTICES.md.
\end{itemize}
\paragraph{Preparation and review.}
The original source and UnsolvedMath contributors retain attribution for the source extract. Codex (GPT-6; exact serving revision not exposed) assisted with definitions, English summaries, scope notes and formatting. No new human mathematical authorship or independently verified solution is asserted. Permanent source ID 30002239 and code OWR-12177-008 checked against the complete supplied registry and current repository. Distinct upstream records, including repeated questions and broad or damaged extracts, are retained for complete coverage of the source CS category. A repeated record is not a new independent conjecture. Complete mathematical-equivalence and current-literature review has not been performed. Preparation and final review follow contribution prompts C and D, with source gaps disclosed. Structural and rendering checks do not establish a new theorem or open status.
\label{end:30002239}
\end{document}
